The measured sketch gives \(h=6.00\,\mathrm{mm}\) and \(b=9.20\,\mathrm{mm}\). The inner and outer lengths are 90 and 103 mm, respectively. Following the course-note estimate, I place the effective root midway between them, giving \(L=96.5\,\mathrm{mm}\). The material inputs are \(E=70.0\,\mathrm{GPa}\) and \(\rho=2700\,\mathrm{kg/m^3}\). Thus,
\[
A=bh=55.2\,\mathrm{mm^2},\quad m=\rho A=0.149\,\mathrm{kg/m},\quad
I=bh^3/12=1.656\times10^{-10}\,\mathrm{m^4},
\]
\[
EI=11.59\,\mathrm{N\,m^2},\qquad a=(EI/m)^{1/4}=2.97\,\mathrm{m/\sqrt{s}}.
\]
Using \(z_1=1.875104\), the initial frequency is
\[
 f_1=\frac{a^2 z_1^2}{2\pi L^2}=530\,\mathrm{Hz},\qquad
 \frac{f_1-512}{512}\,100=3.51\%.
\]
The estimate is above the specified frequency. Since frequency varies as \(L^{-2}\), increasing the estimated length decreases the predicted frequency. Manual trials of 98.0 and 98.2 mm give approximately 514 and 512 Hz. The same adjustment follows directly from the inverse-square relation
\[
\frac{f_{1,\mathrm{new}}}{f_1}=\left(\frac{L}{L_{\mathrm{new}}}\right)^2,\qquad
L_{\mathrm{new}}=L\sqrt{\frac{f_1}{512\,\mathrm{Hz}}}=98.1780\,\mathrm{mm}.
\]
Substitution into the frequency expression gives 512 Hz before rounding.

The initial length is an estimate of the location where each tine begins to bend. Adjusting it by approximately 1.68 mm makes the cantilever estimate agree with the nominal fundamental and provides an effective root location for this model. The three higher predicted bending-mode frequencies, calculated with \(f_j=512(z_j/z_1)^2\), are
\[
 f_2=3.21\,\mathrm{kHz},\qquad f_3=8.98\,\mathrm{kHz},\qquad f_4=17.6\,\mathrm{kHz}.
\]
Their ratios are not integer multiples of the fundamental. The model and measured-sketch interpretation follow the flexural vibration notes, section 2.6, PDF page 186.
\matlabcaption{MATLAB calculations for B.4.}
\lstinputlisting[style=hwmatlab]{matlab/B4.m}
